\chapter{Conclusiones}
\label{chap:conclusiones}

El desarrollo realizado permite concluir que los objetivos planteados al inicio del trabajo se han cumplido.
El resultado final es una aplicación educativa interactiva, accesible desde la web, que permite estudiar de forma visual el entrenamiento y la inferencia de un árbol de decisión y de una red neuronal multicapa.
Para llegar a ese resultado se partió del estudio conceptual de ambos algoritmos y se trasladaron sus operaciones fundamentales a una implementación propia en Godot, manteniendo una correspondencia directa entre el cálculo interno, la representación gráfica y las explicaciones mostradas al usuario.

El proceso seguido comenzó con la validación de Godot como entorno adecuado para construir una aplicación interactiva exportable a navegador.
Esta primera fase permitió comprobar que era posible representar estructuras dinámicas, modificar elementos durante la ejecución y desplegar el proyecto mediante GitHub Pages.
A partir de esa base se incorporó el árbol de decisión, pasando de una estructura editable manualmente a una construcción guiada por el algoritmo ID3.
La aplicación permite observar cómo se crean los nodos, cómo se seleccionan los atributos mediante la métrica basada en entropía y cómo se forman las hojas.
La ventana informativa, las gráficas y la navegación por pasos convierten ese proceso en una explicación progresiva, en la que cada decisión del algoritmo queda asociada a los datos que la provocan.

El segundo bloque principal fue la red neuronal.
Su desarrollo exigió representar una arquitectura formada por capas, neuronas y conexiones densas, y vincular esa representación con un modelo lógico que conserva pesos, sesgos, funciones de activación y parámetros de entrenamiento.
La aplicación permite configurar la red según el conjunto de datos, ejecutar el forward pass paso a paso, mostrar el cálculo de cada neurona, comparar la salida con el valor esperado y representar el backward pass mediante la actualización visible de los parámetros.
De esta forma, el entrenamiento deja de presentarse como una operación opaca y pasa a mostrarse como una secuencia de predicciones, errores y correcciones sobre el modelo.

La fase final del desarrollo consolidó la aplicación como recurso educativo.
Se revisó la navegación general, se adaptaron las vistas a distintos tamaños de pantalla, se mejoró la legibilidad visual de los algoritmos, se integró el renderizado de fórmulas matemáticas y se reforzaron las validaciones de datos.
También se completaron los modos de evaluación tanto en el árbol, donde los datos recorren las ramas aprendidas hasta llegar a una hoja, como en la red, donde la inferencia reutiliza el forward pass sin modificar los parámetros entrenados.
Estas funciones permiten observar no solo cómo se obtiene un modelo, sino también cómo se utiliza después para clasificar o predecir nuevos ejemplos.

El trabajo ha permitido adquirir competencias directamente relacionadas con los objetivos propuestos.
La implementación desde cero de los algoritmos obligó a comprender sus parámetros, sus estructuras internas y las condiciones que aparecen durante su ejecución.
La construcción de la interfaz exigió organizar Escenas, Señales, paneles, animaciones y elementos informativos para que la interacción fuera comprensible.
El despliegue web, junto con la carga y descarga de archivos en navegador, amplió el alcance técnico del proyecto y permitió que la aplicación pudiera consultarse como recurso público de aprendizaje.

En conjunto, el proyecto cumple el propósito principal de ofrecer una herramienta que ayude a comprender paso a paso el funcionamiento interno de dos modelos básicos de aprendizaje automático.
La aplicación final no se limita a mostrar el resultado de un entrenamiento, sino que hace visible el proceso que conduce a ese resultado, acompaña los cálculos con texto, fórmulas y gráficas, y permite experimentar con datos y configuraciones desde una interfaz accesible en web.
